% Источник решений: «семинар 3 (1).pdf», страницы 1--3.
\section{Формулы теории вероятностей}

\noindent Условия:
\href{run:conditions/sem_03.pdf}{семинар 3 (PDF)}.

\subsection{Задача 1}

\begin{solution}
    \textbf{Модель: классическая. Условная вероятность.}
    \[
        \Omega=\{1,\ldots,6\}^3,\qquad \omega=(i_1,i_2,i_3),\qquad
        \mathcal{A}=\{\text{хотя бы одна единица}\}.
    \]
    \textbf{а)} $\mathcal{B}_1=\{i_1,i_2,i_3\text{ попарно различны}\}$.
    \[
        |\mathcal{B}_1|=6\cdot5\cdot4,\qquad
        |\mathcal{A}\cap\mathcal{B}_1|=3\cdot5\cdot4.
    \]
    \[
        \Pr(\mathcal{A}\mid\mathcal{B}_1)
        =\frac{|\mathcal{A}\cap\mathcal{B}_1|}{|\mathcal{B}_1|}
        =\frac{3\cdot5\cdot4}{6\cdot5\cdot4}=\boxed{\frac12}.
    \]
    \textbf{б)} $\mathcal{B}_2=\{i_1=i_3\}$.
    \[
        |\mathcal{B}_2|=6\cdot6=36,\qquad
        |\mathcal{A}\cap\mathcal{B}_2|=5\cdot1+1\cdot6=11.
    \]
    \[
        \Pr(\mathcal{A}\mid\mathcal{B}_2)
        =\frac{|\mathcal{A}\cap\mathcal{B}_2|}{|\mathcal{B}_2|}
        =\boxed{\frac{11}{36}}.
    \]
\end{solution}

\subsection{Задача 2}

\begin{solution}
    \textbf{Модель: классическая, упорядоченная выборка с возвращением.}
    Нумеруем шары; первые $M$ --- белые.
    \[
        \omega=(i_1,\ldots,i_n),\qquad
        \Omega=\{1,\ldots,N\}^n,\qquad |\Omega|=N^n.
    \]
    $\mathcal{A}_k$ --- $k$-й шар белый;
    $\mathcal{B}_m$ --- всего вынуто $m$ белых шаров.
    \[
        |\mathcal{B}_m|=C_n^m M^m(N-M)^{n-m},
    \]
    \[
        |\mathcal{A}_k\cap\mathcal{B}_m|
        =C_{n-1}^{m-1}M^m(N-M)^{n-m}.
    \]
    При $\Pr(\mathcal{B}_m)>0$ и $1\leq m\leq n$:
    \[
        \Pr(\mathcal{A}_k\mid\mathcal{B}_m)
        =\frac{|\mathcal{A}_k\cap\mathcal{B}_m|}{|\mathcal{B}_m|}
        =\frac{C_{n-1}^{m-1}}{C_n^m}
        =\boxed{\frac mn}.
    \]
    При $m=0$ и $\Pr(\mathcal{B}_0)>0$ условная вероятность равна $0$.
\end{solution}

\clearpage
\subsection{Задача 3}

\begin{solution}
    \textbf{Модель: классическая на каждом этапе.
    Формула полной вероятности.}

    $\mathcal{A}$ --- из второй урны вынут белый шар.
    Гипотезы о переложенных шарах:
    \[
        \begin{array}{c|c|c}
            & \Pr(\mathcal{B}_i) & \Pr(\mathcal{A}\mid\mathcal{B}_i)\\[4pt] \hline
            \mathcal{B}_1:\ \text{два белых}
                & \frac35\cdot\frac24=\frac3{10} & \frac6{10}\\[10pt]
            \mathcal{B}_2:\ \text{два чёрных}
                & \frac25\cdot\frac14=\frac1{10} & \frac4{10}\\[10pt]
            \mathcal{B}_3:\ \text{один белый, один чёрный}
                & 2\cdot\frac35\cdot\frac24=\frac6{10} & \frac5{10}
        \end{array}
    \]
    \[
        \begin{aligned}
            \Pr(\mathcal{A})
            &=\Sum{i=1}{3}
                \Pr(\mathcal{A}\mid\mathcal{B}_i)\Pr(\mathcal{B}_i)\\
            &=\frac6{10}\cdot\frac3{10}
             +\frac4{10}\cdot\frac1{10}
             +\frac5{10}\cdot\frac6{10}
             =\boxed{0{,}52}.
        \end{aligned}
    \]
\end{solution}

\subsection{Задача 4}

\begin{solution}
    \textbf{Модель: дискретная. Формула Байеса.}

    $\mathcal{A}$ --- первым стрелял A;
    $\mathcal{B}$ --- первым стрелял B;
    $\mathcal{C}$ --- первый выстрел попал в цель.
    \[
        \Pr(\mathcal{A})=\Pr(\mathcal{B})=\frac12,\qquad
        \Pr(\mathcal{C}\mid\mathcal{A})=\frac12,\qquad
        \Pr(\mathcal{C}\mid\mathcal{B})=\frac45.
    \]
    \[
        \begin{aligned}
            \Pr(\mathcal{A}\mid\mathcal{C})
            &=\frac{\Pr(\mathcal{C}\mid\mathcal{A})\Pr(\mathcal{A})}
            {\Pr(\mathcal{C}\mid\mathcal{A})\Pr(\mathcal{A})
             +\Pr(\mathcal{C}\mid\mathcal{B})\Pr(\mathcal{B})}\\
            &=\frac{\frac12\cdot\frac12}
                    {\frac12\cdot\frac12+\frac45\cdot\frac12}
              =\boxed{\frac5{13}}.
        \end{aligned}
    \]
\end{solution}

\subsection{Задача 5}

\begin{solution}
    \textbf{Модель: классическая, выборка без возвращения.
    Формула умножения.}

    $\mathcal{B}_1$ --- первый шар белый;
    $\mathcal{B}_2$ --- второй шар белый;
    $\mathcal{B}_3$ --- третий шар красный.
    \[
        \begin{aligned}
            \Pr(\mathcal{B}_1\cap\mathcal{B}_2\cap\mathcal{B}_3)
            &=\Pr(\mathcal{B}_1)
              \Pr(\mathcal{B}_2\mid\mathcal{B}_1)
              \Pr(\mathcal{B}_3\mid\mathcal{B}_1\cap\mathcal{B}_2)\\
            &=\frac5{12}\cdot\frac4{11}\cdot\frac7{10}
              =\frac{140}{12\cdot11\cdot10}
              =\boxed{\frac7{66}}.
        \end{aligned}
    \]
\end{solution}
